\section{Controller design on the unstable saddle}
\label{sec:design}

Two architectures are compared at the saddle: a single loop $T\rightarrow q_j$ (SISO) and a cascade $T\rightarrow T_{j,sp}\rightarrow q_j$, with and without feedforward from the coolant supply temperature. All controllers are discrete PI controllers with conditional-integration anti-windup acting on the valve model of Section~\ref{sec:model}, and all gains are computed by the campaign script from the Jacobian of Eq.~\eqref{eq:jacobian}; none is adjusted by hand. A reaction-curve test, the usual starting point of SIMC tuning \cite{skogestad2003simple}, is not available for the single loop because the open-loop response to a step in $q_j$ diverges.

\subsection{Single loop: SIMC for an unstable reduced model}
\label{subsec:siso_design}
At the saddle the transfer function of Eq.~\eqref{eq:gs} has one unstable pole, $1/\tau_u$, two stable poles with time constants $\tau_a>\tau_b$, and a zero with time constant $T_0=-1/a_{11}=\Tzero$~min. The zero is merged with the slower stable lag through $(T_0s+1)/(\tau_as+1)\approx 1/((\tau_a-T_0)s+1)$, which requires $\tau_a>T_0$ and is checked by the script. Skogestad's half rule assigns half of the largest neglected lag to the dominant time constant \cite{skogestad2003simple}; here the dominant pole is unstable and cannot absorb a stable lag, so the remaining lags and half a sampling interval are all lumped into the effective delay. The result is
\begin{equation}
G(s)\approx \frac{k\,e^{-\theta s}}{\tau_us-1},\qquad k=-G(0)=\kU~\frac{\text{K}}{\text{L/min}},
\label{eq:unstable_model}
\end{equation}
with $\tau_u=\tauUmodel$~min and $\theta=\thetaU$~min. The ratio $\theta/\tau_u$ is close to one half, far from the lag-dominant regime, which already signals that this loop is difficult.

Above the frequency $1/\tau_u$ the unstable lag behaves as an integrator with gain $k/\tau_u$, and the SIMC rule for integrating processes \cite{skogestad2003simple} gives
\begin{equation}
K_c=\frac{\tau_u}{k\,(\tau_c+\theta)},\qquad \tau_I=4\,(\tau_c+\theta),
\label{eq:simc_unstable}
\end{equation}
where $\tau_c$ is the desired closed-loop time constant.

\paragraph{Proportional action and the bound $\tau_c<\tau_u-\theta$.}
Let $p(s)=\det(s\mathbf{I}-\mathbf{A}+K_c\mathbf{b}\mathbf{c}^\top)$, with $\mathbf{c}^\top=[0\;1\;0]$, be the characteristic polynomial under proportional feedback. Its constant term is
\begin{equation}
p(0)=\det(-\mathbf{A})\,\bigl(1+K_cG(0)\bigr)=\det(-\mathbf{A})\,\bigl(1-K_ck\bigr).
\end{equation}
At a saddle $\det(-\mathbf{A})<0$, because $\mathbf{A}$ has one positive and two negative real eigenvalues. A stable cubic has all coefficients positive, so stability requires $K_ck>1$: the controller must at least cancel the static loop gain to move the unstable pole into the left half plane. With Eq.~\eqref{eq:simc_unstable}, $K_ck=\tau_u/(\tau_c+\theta)$, and the condition is exactly
\begin{equation}
\tau_c<\tau_u-\theta=\boundU~\text{min}.
\label{eq:tau_bound}
\end{equation}
Equation~\eqref{eq:tau_bound} is necessary for any gain of the form \eqref{eq:simc_unstable}. For this reactor it is also sufficient for proportional control: the column ``P'' of Table~\ref{tab:simc_sweep} is negative for every $\tau_c$ below the bound, and the script verifies that the spectral abscissa of the proportional loop is zero at $K_ck=1$. Note that $k$ in Eq.~\eqref{eq:unstable_model} was defined as $-G(0)$ of the full third-order model, so the bound inherits no error from the model reduction.

\paragraph{Integral action.}
The integral time of Eq.~\eqref{eq:simc_unstable} changes this picture. Table~\ref{tab:simc_sweep} gives the spectral abscissa of the linearized PI loop: it crosses zero at $\tau_c=\tcMaxPI$~min, less than half of the bound \eqref{eq:tau_bound}. In particular the default SIMC choice $\tau_c=\theta$ (last row) gives an unstable loop. We attribute this to the phase lag of the integral term at the frequencies where the loop must overcome the unstable pole; Eq.~\eqref{eq:tau_bound} only accounts for the static gain. For the comparisons below we use $\tau_c=\tauCsiso$~min, giving $K_c=\KcSiso$~L/(min\,K), $\tau_I=\TiSiso$~min and a linear closed-loop spectral abscissa of $\absSiso$~min$^{-1}$. The slowest closed-loop mode therefore decays with time constant $1/|\absSiso|$~min: the single loop can hold the saddle, but only slowly.

\begin{table*}[t]
\centering
\footnotesize
\caption{SIMC PI for the unstable reduced model at the saddle: spectral abscissa of the linearized closed loop for proportional-only (P) and PI action. The last row is the default SIMC choice $\tau_c=\theta$.}
\label{tab:simc_sweep}
\resizebox{\ifdim\width>\textwidth\textwidth\else\width\fi}{!}{%
\begin{tabular}{lccccc}
\toprule
$\tau_c$ (min) & $K_c$ (L/(min K)) & $\tau_I$ (min) & $|K_c k|$ & $\max\Re\lambda$, P & $\max\Re\lambda$, PI \\
\midrule
0.050 & -4.365 & 2.793 & 1.926 & -0.413 & -0.083 \\
0.100 & -4.073 & 2.993 & 1.797 & -0.389 & -0.064 \\
0.200 & -3.593 & 3.393 & 1.585 & -0.351 & -0.029 \\
0.300 & -3.214 & 3.793 & 1.418 & -0.321 & 0.002 \\
0.500 & -2.654 & 4.593 & 1.171 & -0.277 & 0.055 \\
0.648 & -2.351 & 5.185 & 1.038 & -0.059 & 0.087 \\
\bottomrule
\end{tabular}%
}
\end{table*}


\subsection{Cascade: the exact closed-inner-loop sub-block}
\label{subsec:cascade_design}
\paragraph{Slave loop.} The jacket row of Eq.~\eqref{eq:jacobian} gives a first-order model from $q_j$ to $T_j$ with
\begin{equation}
K_2=-\frac{b_3}{a_{33}}=\Ktwo~\frac{\text{K}}{\text{L/min}},\quad \tau_2=-\frac{1}{a_{33}}=\tauTwo~\text{min};
\end{equation}
one sampling interval is taken as delay. Following the usual inner-loop-first sequential design \cite{seborg2010process}, the slave SIMC PI is tuned with $\tau_{c2}=\tau_2$, which gives $K_{c2}=\KcSlave$~L/(min\,K) and $\tau_{I2}=\TiSlave$~min.

\paragraph{What the master loop sees.} If the slave loop held $T_j$ exactly at its setpoint, the jacket state would become an input and the reactor would reduce to the upper $2\times2$ block of Eq.~\eqref{eq:jacobian},
\begin{equation}
\mathbf{A}_{12}=\begin{bmatrix} a_{11} & a_{12}\\ a_{21} & a_{22}\end{bmatrix},\qquad
G_1(s)=\frac{a_{23}\,(s-a_{11})}{\det(s\mathbf{I}-\mathbf{A}_{12})},
\label{eq:a12}
\end{equation}
with static gain
\begin{equation}
K_1=G_1(0)=-\frac{a_{11}\,a_{23}}{\det\mathbf{A}_{12}}.
\label{eq:k1_exact}
\end{equation}
At the saddle the thermal element $a_{22}$ is positive---the reactor heats itself faster than its feed and wall cool it at fixed $T_j$---but the concentration feedback through $a_{12}a_{21}$ dominates: $\operatorname{tr}\mathbf{A}_{12}=\traceSad$~min$^{-1}$ and $\det\mathbf{A}_{12}=\detSad$~min$^{-2}$, so both eigenvalues of $\mathbf{A}_{12}$ lie in the left half plane (largest real part $\absASad$~min$^{-1}$). The instability of the saddle therefore lives in the coupling between reactor and jacket: when the reactor warms, the jacket warms with it and removes less heat. Holding the jacket temperature cuts this positive feedback. With the actual slave loop closed and the master loop open, the linearized system is already stable, with spectral abscissa $\absInner$~min$^{-1}$. The master loop of the cascade does not stabilize anything; it removes offset on a stable plant.

\paragraph{Gain overestimation.} A shortcut in sequential cascade tuning is to estimate the master gain as the ratio of the open-loop gains, $G(0)/K_2$ \cite{seborg2010process,shinskey1996process}. Table~\ref{tab:inner_gain} compares this ratio with Eq.~\eqref{eq:k1_exact}. At the stable nominal point the ratio overestimates the master gain by $\gainErrNom\%$, which only makes a ratio-based master loop more conservative. At the saddle the ratio is $\ratioSad$~K/K against an exact $\KoneSad$~K/K: the sign is wrong. The static gain $G(0)$ of an unstable plant describes an equilibrium that cannot be reached in open loop, and the ratio inherits its sign. A master controller tuned with the ratio would apply positive feedback on the reactor temperature.

\begin{table*}[t]
\centering
\footnotesize
\caption{Static gain seen by the master loop: open-loop ratio approximation versus the exact closed-inner-loop sub-block $\mathbf{A}_{12}$ (Eq.~\eqref{eq:k1_exact}).}
\label{tab:inner_gain}
\resizebox{\ifdim\width>\textwidth\textwidth\else\width\fi}{!}{%
\begin{tabular}{lcccccc}
\toprule
Point & $G(0)$ (K/(L/min)) & $K_2$ (K/(L/min)) & $G(0)/K_2$ (K/K) & $K_{1}$ exact (K/K) & Error (\%) & $\max\Re\lambda(\mathbf{A}_{12})$ \\
\midrule
Nominal & -0.1957 & -0.2416 & 0.810 & 0.576 & 40.5 & -2.433 \\
Saddle & 0.4413 & -0.1864 & -2.367 & 4.521 & -152.4 & -0.340 \\
\bottomrule
\end{tabular}%
}
\end{table*}


\paragraph{Master loop.} Because the closed inner loop is stable, the master loop can be tuned from a reaction curve again. A unit step in $T_{j,sp}$ is applied to the linearized system with the slave loop closed and a first-order-plus-delay model is fitted to the response, giving $K_1=\KoneFit$~K/K (close to the exact value of Eq.~\eqref{eq:k1_exact}), $\tau_1=\tauOneFit$~min and $\theta_1=\thetaOneFit$~min. The master SIMC PI with the default $\tau_{c1}=\theta_1$ gives $K_{c1}=\KcMaster$ and $\tau_{I1}=\TiMaster$~min. The linearized cascade has spectral abscissa $\absCasc$~min$^{-1}$, faster than the single loop, although the master loop uses the default, not a tightened, SIMC setting.

\subsection{Static utility feedforward}
\label{subsec:ff}
A change in the coolant supply temperature is a measured disturbance of the jacket balance. Holding the reactor and jacket at their setpoints, the steady jacket balance gives the flow that keeps the wall heat flux unchanged:
\begin{equation}
q_{ff}(T_{jf}) = \frac{UA_0}{\rho_jc_{pj}}\,\frac{\bar T-\bar T_j}{\bar T_j-T_{jf}}.
\label{eq:ff}
\end{equation}
It is applied as the bias of the slave controller, so the slave integral only has to correct the model error. Equation~\eqref{eq:ff} uses the clean heat-transfer capacity, i.e.\ it carries no information on fouling, and it does not act on unmeasured disturbances such as the feed temperature. For a $+10$~K shift in $T_{jf}$ it prescribes $\ffSad$~L/min, equal to the flow of the new equilibrium ($\qjShockTjf$~L/min); the feedforward is exact in steady state for this disturbance.
